%%%%%%
%
% $Autor: Wings $
% $Datum: 2019-03-05 08:03:15Z $
% $Pfad:  $
% $Version: 4250 $
% !TeX spellcheck = en_GB/de_DE
% !TeX encoding = utf8
% !TeX root = ArduinoIDE2Mac
% !TeX TXS-program:bibliography = txs:///biber
%
%%%%%%


\subsection{Installation (MacOS)}

To install the Arduino IDE, begin by downloading the latest version from the Arduino website Arduino Software \HREF{https://www.arduino.cc/en/software}{Arduino-Software}. Select the version compatible with the specific operating system in use. In this case, Arduino 2.3.2 is being installed for macOS (Sonoma 14.4.1). The setup file is named \HREF{https://www.arduino.cc/en/software}{Arduino-Software} and has a size of 193,600 KB. This file is in Zip format. For those downloading Arduino IDE 2.3.X with Safari, the file will automatically extract upon download completion, while other browsers may require manual extraction. The figure below displays the latest offline version of Arduino IDE 2.3.3 ~\ref{fig:ArduinoIDECreateAgentInstallation}, which is also compatible with all operating systems.




\begin{center}
	
	\includegraphics[width=\linewidth]{Arduino/ArdiunoIDE2/macDownload.png}
	%\label{fig:ArduinoIDE Create Agent Installation}
	\captionof{figure}{ArduinoIDE Create Agent Installation}
	\label{fig:ArduinoIDECreateAgentInstallation}
	
	
	
	
	
	
	\Mynote{Die Installation muss aktualisiert werden, die folgenden Bilder müssen mit einem Mac ebenso aktualisiert werden}
	
\end{center}

\begin{center}
	\includegraphics[width=0.7\linewidth]{Arduino/ArdiunoIDE2/OpenTheDownloadFolder.png}
	\captionof{figure}{Open the Download folder}
\end{center}

Copy the Arduino application bundle into the application's folder (or elsewhere on the
computer) then it looks like Figure   ~\ref{fig:CopytotheAppliationsFolder}.

\begin{center}
	
	\includegraphics[width=0.7\linewidth]{Arduino/ArdiunoIDE2/CopyToTheApplicationsFolder.png}
	\captionof{figure}{Copy to the Applications Folder}
	\label{fig:CopytotheAppliationsFolder}
\end{center}

It can be seen from the Figure ~\ref{fig:ArduinoIDESketch2} that the basic Arduino sketch has two parts. 

\begin{itemize}
	
	\item \PYTHON{void setup()}: This function returns void
	and performs initializations such as setting the output LED color and specifying the core.
	\item \PYTHON{void loop()}: In this function, specific operations to be executed within the loop are defined. The code for each operation is enclosed within curly braces \PYTHON{\{\}}, and each function has a return type. In this case, the return type is \menu{void}.
	
\end{itemize}

\begin{center}
	
	
	
	\includegraphics[width=0.7\linewidth]{Arduino/ArdiunoIDE2/ArduinoIDESketch2.png}
	\captionof{figure}{Arduino IDE Sketch}
	\label{fig:ArduinoIDESketch2}
	
	
	
	\Mynote{Figure 1.22 muss größer und alle Bilder in diesem Kapitel ohne Schwarzmodus}	
	
\end{center}

